%% Resumo
\begin{resumo}
Atualmente, falhas no cálculo da execução penal têm contribuído para a permanência indevida de apenados no sistema prisional brasileiro, agravando um cenário já marcado pelo aumento da população carcerária. Entre os fatores envolvidos, destaca-se um dos mais recorrentes: erros no Processo de Execução Criminal (PEC). A PEC é o conjunto de procedimentos que acompanha o cumprimento da pena, incluindo cálculos, concessão de benefícios e registro das ações de trabalho e/ou estudo para remição, conforme a Lei de Execução Penal (LEP). Quando realizado manualmente, esse cálculo está suscetível a erros que passam despercebidos, podendo resultar em períodos de prisão mais longos do que o legalmente previsto ou em solturas antecipadas indevidas, afetando a gestão de vagas e a segurança pública. O cumprimento correto da pena e a efetivação do cálculo, considerando as remições obtidas pelo apenado, contribuem para que o indivíduo avance no processo de ressocialização, reduzindo a reincidência e colaborando para a segurança da sociedade. Dessa forma, o objetivo deste projeto é automatizar o cálculo da execução penal, reduzindo falhas humanas, garantindo maior precisão e agilidade no cumprimento das penas, utilizando
ferramentas que possibilitam a extração de dados referentes ao apenado durante seu
tempo de comprimento de pena, como o número de dias trabalhados e/ou estudados, para que o sistema possa calcular automaticamente a remição e o tempo restante da pena, proporcionando uma gestão mais eficiente do sistema prisional e contribuindo para a ressocialização dos apenados.
\vspace{\onelineskip}
\noindent

\textbf{Palavras-chaves}: execução penal; sistema prisional; automação; cálculo automático; Lei de Execução Penal; PEC; extração de dados.
\end{resumo}